\section{Authority must survive every transformation}\label{sec:C17:authority}
Authority is information about who may direct which effects under what conditions. It should not be inferred from memory frequency, semantic relevance, or the confidence of a generated summary. A retrieved page can be highly relevant to a task while being untrusted as an instruction source. A stored skill can be useful while lacking authority to change the task. A prior session can contain valid historical permission that no longer applies.

Progent represents allowed tool calls through symbolic policies and distinguishes narrowing from expansion~\cite{corpus250411703}. Under complete mediation, an unapproved update can be restricted to a subset of the previous allowed set. Write this as $\mathcal A_{t+1}\subseteq\mathcal A_t$ for automatically accepted narrowing. The comparison must use the same current authority epoch; an allowed set restored from an old snapshot is not the current grant. The guarantee depends on correct policy interpretation, a trusted enforcement boundary, and every consequential call passing through it. It does not guarantee that all allowed calls are useful or semantically correct.

Compaction can preserve authority in two ways. It can retain the original authoritative statement verbatim, or it can preserve a validated typed representation whose semantics are enforced independently. A generated paraphrase without validation can weaken or broaden a condition. For example, ``prepare a draft for review'' can become ``complete the publication task'' after repeated summarization. The model may then believe a broader objective even if the external monitor still blocks publishing. Preserving both behavioral guidance and enforcement reduces this mismatch.

Concurrent updates create races. A worker may choose an action under authority version $a$ while a user narrows scope to $a+1$. Checking permission only at worker creation is insufficient if the effect occurs later. The dispatcher should evaluate the current applicable policy or require a validated authorization lease whose semantics cover that interval. Canceling the worker is not equivalent to preventing an already dispatched effect.

The proposed test records authority at task creation, context assembly, action proposal, and effect dispatch. Inject a narrowing correction between each pair of stages. Check whether stale proposals are rejected and whether the user-visible state accurately reports any in-flight uncertainty. This targets authority consistency directly instead of relying on a general prompt-injection benchmark to discover the race accidentally.

Authority preservation also constrains learning. A memory writer should not turn a successful prohibited action into a reusable recommendation merely because the environment rewarded it. Outcome feedback and authorization are separate channels. A system can perform well on a task-success metric while learning procedures that violate its intended scope. The evaluation must retain both dimensions and report their interaction.

